% 근거: chapter4-prep-v1/A4/denominators.json,
%       chapter4-prep-v1/A6/observation-frequencies.jsonl,
%       chapter4-prep-v1/A6/proposed-three-candidates.json,
%       chapter4-prep-v1/A6/casebook.jsonl
\subsection{오류 유형 및 관측 사례}
\label{sec:ch4_errors}

오류 분석에서는 exact gold의 회수ㆍ선택 상태와 법률적 정오 판단을
구분하였다. 문항별 고유 gold 출현 443건을 분모로 최초 탈락 단계를
세면, corpus 부재는 33건, corpus에는 있지만 공통 BM25 후보에 없는
경우는 75건으로 두 조건에서 같았다. BM25 후보에 있으나 shortlist에서
탈락한 gold는 수정 CE 42건, Jev 13건이었고, shortlist에는 있으나
최종 선택되지 않은 gold는 각각 60건과 72건이었다.
최종 잔존은 233건과 250건으로, 각 조건의 최초 탈락과 잔존을 합하면
443건이 된다. 이는 전역 고유 gold 386개와 다른 출현 단위 집계다.

Jev는 shortlist 탈락을 줄였지만 후속 미선택도 남겼다.
이 미선택에는 후보 목록당 하나를 고르는 제약과 중복 선택이 작용할 수
있으므로, 관측된 단계를 곧바로 LLM의 법률 판단 오류로 명명하지 않았다.
gold 외 근거가 추가된 문항도 수정 CE 126개, Jev 116개였으나,
해당 근거의 법률적 무관성을 뜻하지는 않는다. 문항 수준의 관측 유형은
서로 겹칠 수 있으므로 유형별 빈도를 합해 전체 오류 문항 수로 사용하지 않는다.

대표 사례는 Jev 개선, 수정 CE 우세, 지표 불일치 묶음에서
근거 F1 차이 절댓값의 중앙값에 가장 가까운 문항을 선정하고,
동률은 query\_id 사전순으로 정하였다.
표~\ref{tab:ch4_cases}의 세 사례는 자동 관측을 설명하기 위한 자료이며,
법률 전문가의 판정이 완료된 사례가 아니다.

\begin{table}[htbp]
    \centering
    \caption{대표 관측 사례의 자동 점수 변화 (수정 CE$\to$Jev)}
    \label{tab:ch4_cases}
    \small
    \setlength{\tabcolsep}{4pt}
    \begin{tabular}{@{}llrr@{}}
        \hline
        사례 & 질문 주제 & 근거 F1 & LF-Eval \\
        \hline
        A & 식품 보고 & $50.00\to100.00$ & $65.29\to99.18$ \\
        B & 복수 청구 & $80.00\to40.00$  & $71.79\to52.67$ \\
        C & 사면 효과 & $40.00\to0.00$   & $31.10\to30.55$ \\
        \hline
    \end{tabular}
    \par\smallskip
    \begin{minipage}{\linewidth}
        \footnotesize
        점수는 0--100 척도이다. 사례 식별자는 다음과 같다.\\
        A: \texttt{qa\_124\_2hop\_480\_rand}\\
        B: \texttt{qa\_324\_2hop\_515}\\
        C: \texttt{qa\_443\_3hop\_703}
    \end{minipage}
\end{table}

사례 A는 식품 생산실적과 위해식품 회수계획의 보고 상대방을 묻는
2-hop 문항이다. 수정 CE는 식품위생법 제42조제2항과 생산실적 관련
시행규칙을 선택한 반면, Jev는 gold인 제42조제2항과 제45조제1항을
모두 선택하였다. Jev 답변에는 회수계획의 보고 상대방으로
시장ㆍ군수ㆍ구청장이 포함되었고, 근거 F1과 LF-Eval이 함께 높아졌다.
이는 최종 근거 선택과 답변 문구의 차이가 함께 나타난 관측 사례다.

사례 B는 여러 당사자를 상대로 복수 청구를 제기하는 절차에 관한
2-hop 문항이다. 수정 CE는 gold인 민사소송법 제79조제1항과
제253조를 모두 선택하였다. Jev에서는 제79조제1항이 shortlist에서
탈락하고, 최종 근거에 제83조제1항이 포함되었다.
답변에서도 독립당사자참가 관련 내용과 공동소송참가 관련 내용의
차이가 나타났으며, 이 문항은 두 자동 지표 모두 수정 CE가 높았다.
각 제도의 실제 적용 가능성은 별도 법률 검토 사항으로 남겼다.

사례 C는 사면과 복수 범죄의 처리에 관한 3-hop 문항이다.
gold 3개 중 하나는 corpus에 정확히 매핑되지 않았고 다른 하나는
BM25 후보에 없었으며, 생성된 후보 목록도 2개였다.
수정 CE는 gold 하나를 최종 선택하였지만 Jev는 exact gold를
선택하지 않았다. 반면 답변 Token F1은 18.60점에서 25.88점으로
높아져, 근거 F1과 문자열 기반 답변 지표가 반대 방향을 보였다.
사면 유형과 범죄의 경합 관계, gold 외 선택 조문의 타당성은
전문가 확인 전 확정하지 않았다.

세 사례는 모두 기존 사람 평가 60문항 표본 밖의 별도 정성 검토 대상이며,
블라인드 표본을 교체하거나 사람 평가 결과에 추가하지 않았다.
위 사례 분석은 단계별 탈락과 지표 불일치를 설명하며,
법률적 오류 여부와 제공 근거의 주장 지지는 미판정 상태로 구분한다.
